\section*{Chapitre \ref{ch:dv:greeks-and-the-hedging-pnl} --- Grecques et P\&L de couverture}

\begin{solution}{exo:dv:greeks-and-the-hedging-pnl:1}
Delta 0.0230, gamma 0.1381, \omterm{def:dv:greeks-and-the-hedging-pnl:greeks}{vega} 0.2302 par point, \omterm{def:dv:greeks-and-the-hedging-pnl:greeks}{theta} $-0.0757$ par jour calendaire.
\end{solution}

\begin{solution}{exo:dv:greeks-and-the-hedging-pnl:2}
Avec $r=q=0$~: $\Theta=-\tfrac12\sigma^2S^2\Gamma=-\tfrac12\times0.04\times10\,000\times0.1381
=-27.63$ par an, $-0.0757$ par jour, soit le \omterm{def:dv:greeks-and-the-hedging-pnl:greeks}{theta} analytique.
\end{solution}

\begin{solution}{exo:dv:greeks-and-the-hedging-pnl:3}
P\&L de gamma $5\,000\times0.02^2=2.00$~; \omterm{def:dv:greeks-and-the-hedging-pnl:greeks}{theta} $-5\,000\times0.2^2/365=-0.55$~; net $+1.45$.
\end{solution}

\begin{solution}{exo:dv:greeks-and-the-hedging-pnl:4}
Chaque intervalle rapporte $\tfrac12\Gamma S^2(x^2-\sigma^2\delta t)$, quadratique dans le rendement
$x$. Une action qui dérive vers 120 par petits pas (chacun en deçà du mouvement de point mort) perd
du \omterm{def:dv:greeks-and-the-hedging-pnl:greeks}{theta} chaque jour~: le \omterm{def:dv:greeks-and-the-hedging-pnl:cashgamma}{straddle} long, couvert, perd de l'argent alors que, non couvert, il
aurait touché 20 à l'échéance. La couverture a vendu la dérive au fur et à mesure qu'elle se produisait.
\end{solution}

\begin{solution}{exo:dv:greeks-and-the-hedging-pnl:5}
En quadruplant le nombre de transactions~: l'erreur décroît comme $1/\sqrt N$. Chaque transaction rééquilibre alors une variation de delta
de moitié environ, si bien que la quantité totale échangée, et donc le coût, croissent comme $\sqrt N$~:
deux fois le coût pour moitié moins de bruit (le compromis de Leland).
\end{solution}

\begin{solution}{exo:dv:greeks-and-the-hedging-pnl:6}
$\sqrt{\pi/4}\times23.02\times0.2/\sqrt{21}=0.890$, à 3~\% près de la valeur simulée 0.865.
\end{solution}

\begin{solution}{exo:dv:greeks-and-the-hedging-pnl:7}
Moyenne $-1.147$, écart-type 1.082~: la moyenne est inchangée, et avec des couvertures quotidiennes
le bruit de couverture discrète, et non le choix de la volatilité, domine la dispersion.
\end{solution}

\begin{solution}{exo:dv:greeks-and-the-hedging-pnl:8}
Un \omterm{def:dv:greeks-and-the-hedging-pnl:cashgamma}{straddle} vendu encaisse du \omterm{def:dv:greeks-and-the-hedging-pnl:greeks}{theta} la plupart des jours et paie du gamma les quelques jours de fortes
variations~; lorsque la volatilité réalisée dépasse l'implicite, les pertes quadratiques des jours de grands mouvements dépassent les
gains réguliers. Une proportion élevée de jours gagnants est la signature de la vente
d'optionnalité, non une preuve d'erreur de valorisation~; la perte résulte de l'identité gamma--theta à
l'œuvre.
\end{solution}

\begin{solution}{pb:dv:greeks-and-the-hedging-pnl:1}
\textbf{1.} 4.6059 par \omterm{def:dv:greeks-and-the-hedging-pnl:cashgamma}{straddle}~; 460\,595~USD.
\textbf{2.} Chaque \omterm{def:dv:greeks-and-the-hedging-pnl:cashgamma}{straddle} a un delta de 0.0230~; la position vendue de 1\,000 avec un multiplicateur de 100 correspond à $-2\,303$
actions~; le desk achète 2\,303 actions.
\textbf{3.} \omterm{def:dv:greeks-and-the-hedging-pnl:cashgamma}{Cash gamma} 690.70~; gamma pour un mouvement de 1~\%~: 13.81 (en cash delta).
\textbf{4.} \omterm{def:dv:greeks-and-the-hedging-pnl:greeks}{Vega} de 23\,023~USD par point (position vendue, donc une perte quand la volatilité monte)~; \omterm{def:dv:greeks-and-the-hedging-pnl:greeks}{theta}
de 7\,569~USD par jour au profit du desk.
\textbf{5.} 1.047 sur 365 jours~; 1.260 sur 252.
\textbf{6.} 5.7570.
\textbf{7.} 1.1510 par \omterm{def:dv:greeks-and-the-hedging-pnl:cashgamma}{straddle}~; 115\,104~USD.
\textbf{8.} $5\times0.2302=1.1512$.
\textbf{9.} La valeur du \omterm{def:dv:greeks-and-the-hedging-pnl:cashgamma}{straddle} est presque linéaire en volatilité à la monnaie (le \omterm{def:dv:greeks-and-the-hedging-pnl:greeks}{volga} est
faible), si bien que le \omterm{def:dv:greeks-and-the-hedging-pnl:greeks}{vega} multiplié par l'écart fournit une estimation du premier ordre fiable.
\textbf{10.} Non~: la couverture continue donne la même espérance, sans dispersion
due au calendrier de couverture.
\textbf{11.} Moyenne $-1.145$, écart-type 1.183.
\textbf{12.} 0.620.
\textbf{13.} 87.3~\%.
\textbf{14.} Moyenne $-1.147$, écart-type 1.082.
\textbf{15.} Moyenne 0.003, écart-type 0.865~: pur bruit de couverture discrète.
\textbf{16.} Que l'action ne dégage pas plus de 20 sur le mois, après frais et
avec une marge pour le bruit de couverture.
\textbf{17.} Limiter le \omterm{def:dv:greeks-and-the-hedging-pnl:greeks}{vega} (et le gamma à l'approche de l'échéance) de sorte qu'un choc de volatilité
plausible, soit dix points, et deux écarts-types du bruit de couverture restent dans la
limite de perte mensuelle.
\textbf{18.} Le \omterm{def:dv:greeks-and-the-hedging-pnl:greeks}{vega}, lu avec un scénario de volatilité réalisée, et la
concentration du gamma près du prix d'exercice.
\textbf{19.} Une perte espérée de 115\,100~USD et un écart-type de 118\,300~USD
avec une couverture quotidienne.
\textbf{20.} Le \omterm{def:dv:greeks-and-the-hedging-pnl:cashgamma}{cash gamma} multiplié par l'écart entre la variance réalisée et la variance implicite,
sommé sur la durée de vie de l'option.
\end{solution}

\begin{solution}{iq:dv:greeks-and-the-hedging-pnl:1}
La position a perdu une semaine de \omterm{def:dv:greeks-and-the-hedging-pnl:greeks}{theta}~: l'option a perdu de sa valeur, et aucun P\&L de gamma n'est venu compenser cette perte.
D'après l'identité gamma--theta, une option longue couverte ne rapporte que lorsque le sous-jacent
bouge davantage que ne le suppose la volatilité implicite.
\iqlookfor{le \omterm{def:dv:greeks-and-the-hedging-pnl:greeks}{theta} comme loyer du gamma.}
\end{solution}

\begin{solution}{iq:dv:greeks-and-the-hedging-pnl:2}
$\Theta\approx-\tfrac12\Gamma S^2\sigma^2$~: la décroissance temporelle d'une option couverte est
exactement ce que rapporte son gamma si le sous-jacent bouge d'un écart-type de la
volatilité implicite chaque jour. Être long en gamma revient à payer du \omterm{def:dv:greeks-and-the-hedging-pnl:greeks}{theta} pour la convexité.
\iqlookfor{l'identité, et «~point mort à l'implicite~».}
\end{solution}

\begin{solution}{iq:dv:greeks-and-the-hedging-pnl:3}
Développer l'option en série de Taylor~: $\Theta\delta t+\Delta\delta S+\tfrac12\Gamma\delta S^2$~; la
couverture supprime $\Delta\delta S$~; l'identité remplace $\Theta$ par $-\tfrac12\Gamma S^2
\sigma_{\mathrm{imp}}^2$ (taux nuls)~: P\&L $=\tfrac12\Gamma S^2\bigl((\delta S/S)^2-
\sigma_{\mathrm{imp}}^2\delta t\bigr)$.
\iqlookfor{le développement au second ordre et l'annulation.}
\end{solution}

\begin{solution}{iq:dv:greeks-and-the-hedging-pnl:4}
Si la prévision est juste, le P\&L à l'échéance est figé, mais la valorisation au prix de
marché fluctue fortement au jour le jour~; à l'implicite, le P\&L quotidien s'explique par le gamma
et le \omterm{def:dv:greeks-and-the-hedging-pnl:greeks}{theta}, mais le total dépend de l'endroit où surviennent les grands mouvements par rapport
au prix d'exercice. Couvrir à l'implicite est le choix par défaut, car la prévision est incertaine
et car les rapports de risque et les limites sont exprimés en termes implicites.
\iqlookfor{les deux propriétés, et pourquoi les desks choisissent l'implicite.}
\end{solution}

\begin{solution}{iq:dv:greeks-and-the-hedging-pnl:5}
Monte Carlo~: les réévaluations à la hausse et à la baisse utilisent des nombres aléatoires indépendants, de sorte que leur
différence est du bruit divisé par $h^2$~; utiliser des nombres aléatoires communs, des estimateurs pathwise ou
likelihood-ratio (chapitre 23). Option digitale à l'approche de l'échéance~: le payoff est
discontinu, le gamma est un pic plus étroit que le bump, et la différence mesure
la taille du bump~; présenter le gamma sur un bump correspondant à l'horizon de risque, et gérer
l'option digitale par un call spread (chapitre 15).
\iqlookfor{bruit issu de tirages indépendants, et payoffs non lisses.}
\end{solution}

\begin{solution}{iq:dv:greeks-and-the-hedging-pnl:6}
Le \omterm{def:dv:greeks-and-the-hedging-pnl:greeks}{vega} cumule le gamma sur la durée de vie des options, mais le P\&L pondère la variance réalisée
de chaque jour par le \omterm{def:dv:greeks-and-the-hedging-pnl:cashgamma}{cash gamma} de ce jour. Un portefeuille dont le gamma se situait aux strikes
que le spot a visités lors de ses jours volatils gagne davantage qu'un portefeuille dont le gamma
se trouvait ailleurs~; cet écart peut être de l'ordre de toute l'espérance du P\&L, comme le montre
la dispersion de la \cref{fig:dv:greeks-and-the-hedging-pnl:which}.
\iqlookfor{pondération du gamma dans le temps et dans le spot.}
\end{solution}
